\section*{अध्याय \ref{ch:g6:solutions-and-solubility} --- विलयन और विलेयता}

\begin{solution}{exo:g6:solutions-and-solubility:1}
\omterm{def:g6:solutions-and-solubility:solute}{विलेय}: चीनी। \omterm{def:g6:solutions-and-solubility:solute}{विलायक}: पानी। हाँ, \omterm{def:g6:solutions-and-solubility:solute}{विलायक} पानी है: यह एक
\omterm{def:g6:solutions-and-solubility:solute}{जलीय विलयन} है।
\end{solution}

\begin{solution}{exo:g6:solutions-and-solubility:2}
ऐसा विलयन जो अपने \omterm{def:g6:solutions-and-solubility:solute}{विलेय} की और मात्रा नहीं घोल सकता। उसमें डाला गया \omterm{def:g6:solutions-and-solubility:solute}{विलेय}
बिना घुले तली में रह जाता है।
\end{solution}

\begin{solution}{exo:g6:solutions-and-solubility:3}
$250 + 20 = \qty{270}{g}$।
\end{solution}

\begin{solution}{exo:g6:solutions-and-solubility:4}
एथेनॉल: \omterm{def:g6:solutions-and-solubility:miscible}{मिश्रणीय}। खाना पकाने का तेल: \omterm{def:g6:solutions-and-solubility:miscible}{अमिश्रणीय}। सिरका: \omterm{def:g6:solutions-and-solubility:miscible}{मिश्रणीय} (वह
ज़्यादातर पानी है)।
\end{solution}

\begin{solution}{exo:g6:solutions-and-solubility:5}
प्रोपेनोन (ऐसीटोन)। GHS02: अत्यधिक ज्वलनशील; GHS07: आँखों में जलन करता है,
वाष्प से उनींदापन या चक्कर आ सकता है।
\end{solution}

\begin{solution}{exo:g6:solutions-and-solubility:6}
$360 \times \frac{500}{1000} = \qty{180}{g}$।
\end{solution}

\begin{solution}{exo:g6:solutions-and-solubility:7}
\qty{200}{g} पानी अधिकतम $360 \times \frac{200}{1000} =
\qty{72}{g}$ घोलता है। नहीं: $80 - 72 = \qty{8}{g}$ तली में रह जाता है।
\end{solution}

\begin{solution}{exo:g6:solutions-and-solubility:8}
विलयन का द्रव्यमान $225 + 25 = \qty{250}{g}$ है; $\frac{25}{250} =
\qty{10}{\percent}$।
\end{solution}

\begin{solution}{exo:g6:solutions-and-solubility:9}
पहले बीकर में, जहाँ सारा नमक घुल चुका है और विलयन अभी
संतृप्त नहीं है।
\end{solution}

\begin{solution}{exo:g6:solutions-and-solubility:10}
बंद बोतल में कार्बन डाइऑक्साइड दाब के साथ घोली गई थी। बोतल खोलने पर दाब
घट जाता है: पानी कम गैस रख सकता है, और अतिरिक्त गैस बुलबुलों के रूप में
बाहर आ जाती है।
\end{solution}

\begin{solution}{exo:g6:solutions-and-solubility:11}
$2000 \div 0.04 = 50\,000$ गुना ज़्यादा।
\end{solution}

\begin{solution}{exo:g6:solutions-and-solubility:12}
कोई गैस गरम पानी में कम \omterm{def:g2:mixing-and-dissolving:dissolve}{घुलती} है: गरम तालाब में घुली हुई ऑक्सीजन कम
होती है, और मछलियों को ऑक्सीजन की कमी पड़ती है।
\end{solution}

\begin{solution}{pb:g6:solutions-and-solubility:1}
\textbf{1.} \omterm{def:g6:solutions-and-solubility:solute}{विलेय}: नमक। \omterm{def:g6:solutions-and-solubility:solute}{विलायक}: पानी।

\textbf{2.} पानी में जितना नमक हो सकता है, उतना है: उसमें और नमक नहीं
घुल सकता।

\textbf{3.} तली में पड़ा बिना घुला नमक सिद्ध करता है कि खारा पानी
संतृप्त है: उसमें अधिक से अधिक संभव नमक है।

\textbf{4.} $360 \times 20 = \qty{7200}{g} = \qty{7.2}{kg}$।

\textbf{5.} $20 + 7.2 = \qty{27.2}{kg}$।

\textbf{6.} $\frac{7.2}{27.2} \approx 0.26$: लगभग \qty{26}{\percent}।

\textbf{7.} नहीं: पानी \qty{7.2}{kg} घोल सकता था। वह अब भी
$7.2 - 5 = \qty{2.2}{kg}$ घोल सकता है।

\textbf{8.} $20 - 5 = \qty{15}{kg}$ पानी (\qty{15}{L})।

\textbf{9.} $360 \times 15 = \qty{5400}{g} = \qty{5.4}{kg}$।

\textbf{10.} वह विलयन से ठोस क्रिस्टलों के रूप में बाहर आता है: वह
क्रिस्टल बनकर तली में बैठ जाता है।

\textbf{11.} टंकी की तली में नए सफ़ेद क्रिस्टल प्रकट होते हैं और जमा होते
जाते हैं।

\textbf{12.} $7.2 - 5.4 = \qty{1.8}{kg}$ नमक क्रिस्टल बनकर जमा है।
\end{solution}
